\chapter{Lumbar laminectomy}
\index{Lumbar laminectomy}

\section*{Definition and Overview}
A lumbar laminectomy, also known as decompression surgery, is a surgical procedure performed on the lower back to relieve pressure on the spinal cord or spinal nerve roots. It involves the removal of the lamina---the bony arch on the back of the vertebra---along with any bone spurs or ligaments that are narrowing the spinal canal. This decompression helps alleviate symptoms such as pain, numbness, and weakness in the legs.

\section*{Epidemiology}
Lumbar laminectomy is one of the most common spinal surgeries, particularly among older adults. Spinal stenosis, the primary indication for this procedure, is a degenerative condition that typically affects people over the age of 60. There is a slight male predominance, and the incidence increases with age as degenerative changes in the spine accumulate.

\section*{Aetiology and Pathophysiology}
Lumbar spinal stenosis occurs when the space within the lower spinal canal narrows, compressing the nervous tissue. The laminectomy surgically addresses this mechanical compression:

\begin{itemize}
    \item \textbf{Spinal stenosis:} Degenerative changes, such as osteoarthritis of the facet joints, ligamentum flavum hypertrophy, and bulging discs, reduce the canal diameter.
    \item \textbf{Bony compression:} Osteophytes (bone spurs) project into the canal or neural foramina, compressing the cauda equina or exiting nerve roots.
    \item \textbf{Surgical decompression:} Removing the lamina and associated hypertrophied ligaments increases the spinal canal space, restoring blood flow to the nerves and relieving mechanical pressure.
\end{itemize}

\section*{Clinical Features}
Patients requiring or recovering from a lumbar laminectomy exhibit characteristic neurological and musculoskeletal symptoms:

\begin{itemize}
    \item \textbf{Neurogenic claudication:} Leg pain, numbness, or cramping that occurs when walking or standing and is relieved by sitting or bending forward.
    \item \textbf{Radiculopathy:} Sharp, shooting pain (sciatica) radiating from the lower back down one or both legs, following a specific nerve root pathway.
    \item \textbf{Sensory and motor deficits:} Weakness in the legs, difficulty lifting the foot (foot drop), or numbness in the legs and buttocks.
    \item \textbf{Post-operative surgical pain:} Localized pain and stiffness at the incision site in the lower back during the initial recovery phase.
\end{itemize}

\section*{Differential Diagnosis}
Other conditions causing lower back and leg pain must be distinguished from spinal stenosis:

\begin{itemize}
    \item \textbf{Vascular claudication:} Leg pain caused by peripheral arterial disease, which is relieved by standing still and is unaffected by spinal posture.
    \item \textbf{Hip osteoarthritis:} Joint pain that can mimic radicular pain and is exacerbated by hip movement.
    \item \textbf{Lumbar disc herniation:} Acute nerve compression, more common in younger adults, often presenting without diffuse stenosis.
    \item \textbf{Diabetic neuropathy:} Symmetric burning pain and numbness in the feet, independent of walking.
\end{itemize}

\section*{When to Seek Medical Care}
Patients with chronic lower back and leg pain that limits walking or daily activities should consult a GP or spinal specialist. Immediate emergency medical attention is critical if symptoms of cauda equina syndrome develop, as this is a spinal surgery emergency.

\textbf{Contact your local emergency services for:}
\begin{itemize}
    \item Sudden loss of bowel or bladder control (incontinence or retention)
    \item "Saddle anaesthesia" (numbness in the groin, buttocks, or inner thighs)
    \item Progressive or sudden weakness in one or both legs, making standing or walking impossible
    \item Severe, unmanageable back or leg pain accompanied by high fever or chills
    \item Signs of wound infection, such as pus discharge, increasing redness, or severe swelling around the surgical site
\end{itemize}

\section*{Diagnosis and Investigations}
Diagnosis is based on clinical presentation and confirmed with advanced imaging:

\begin{itemize}
    \item \textbf{MRI scan:} The gold standard investigation to visualize the spinal canal, nerve roots, and degree of ligamentous or bony narrowing.
    \item \textbf{CT scan or CT myelogram:} Used if MRI is contraindicated (e.g., in patients with certain pacemakers) to assess bony structures.
    \item \textbf{X-rays of the lumbar spine:} Performed to evaluate spinal alignment, disc space narrowing, and the presence of instability or spondylolisthesis.
    \item \textbf{Electrodiagnostic studies (EMG/NCS):} Done occasionally to confirm nerve compression and rule out peripheral neuropathy.
\end{itemize}

\section*{Management}
Management involves pre-operative preparation, the surgical procedure, and structured post-operative rehabilitation:

\begin{itemize}
    \item \textbf{Conservative measures:} A trial of physiotherapy, analgesics (like NSAIDs), and epidural steroid injections before considering surgery.
    \item \textbf{Decompression surgery:} Under general anaesthesia, a midline incision is made in the lower back, muscles are retracted, and the lamina is removed.
    \item \textbf{Mobilisation:} Encouraging the patient to stand and walk on the day of surgery or the following day to prevent deep vein thrombosis.
    \item \textbf{Pain management:} Using a combination of oral analgesics, muscle relaxants, and local ice therapy post-operatively.
    \item \textbf{Rehabilitation:} Structured physiotherapy starting a few weeks post-surgery to strengthen core muscles, improve flexibility, and guide safe return to normal activities.
\end{itemize}

\section*{Complications}
\begin{itemize}
    \item Dural tear (cerebrospinal fluid leak), which may require intra-operative repair or bed rest
    \item Wound infection or discitis (infection of the disc space)
    \item Epidural haematoma compressing nerve roots, requiring urgent surgical evacuation
    \item Persistent pain (failed back surgery syndrome) or nerve root damage
    \item Spinal instability, which might eventually require spinal fusion surgery
    \item General surgical risks, such as deep vein thrombosis or pulmonary embolism
\end{itemize}

\section*{Prognosis and Outcomes}
The prognosis after lumbar laminectomy is generally favourable, with 70\% to 80\% of patients experiencing significant relief from leg pain and claudication. Recovery times vary; most patients return to light activities within a few weeks and can resume driving and light work within 4 to 6 weeks. Complete recovery and maximum benefit are typically seen 3 to 6 months post-operatively. Back stiffness may persist, and degenerative changes elsewhere in the spine can occur over time.

\section*{Prevention and Self-Care}
\begin{itemize}
    \item Maintain a healthy weight to reduce the mechanical load on the lumbar spine
    \item Engage in regular, low-impact exercise (e.g., walking or swimming) to maintain spinal flexibility and strength
    \item Practice proper lifting techniques: bend the knees and keep the back straight
    \item Avoid prolonged sitting or standing in one position; take frequent movement breaks
    \item Quit smoking, as nicotine impairs spinal bone healing and accelerates disc degeneration
    \item Strictly follow the post-operative restriction against bending, twisting, or lifting heavy objects during early recovery
\end{itemize}

\section*{Sources and Further Reading}
The following organisations provide reliable information on lumbar laminectomy:

\begin{itemize}
    \item Spine Society of Australia. \textit{Patient guides on lumbar spinal stenosis and decompression surgery.} https://www.spinesociety.org.au/
    \item Mayo Clinic. \textit{Laminectomy procedure overview, risks, and recovery.} https://www.mayoclinic.org/
    \item NHS. \textit{Guide to lumbar decompression surgery, including what happens and recovery.} https://www.nhs.uk/
    \item Healthdirect Australia. \textit{Information on lumbar laminectomy and spinal stenosis.} https://www.healthdirect.gov.au/
    \item American Academy of Orthopaedic Surgeons. \textit{Patient education on lumbar laminectomy.} https://orthoinfo.org/
\end{itemize}
